

Continuing with the established notation, let $w_i \in \R$ be the bound on player $i$'s utility, and let $\bmu$ be probability simplices defining the mixed strategies of players. The multilinear feasibility program in $\bw$ and $\bmu$ for finding NE in general-sum games is
\begin{equation}\label{eq:sturmfels}
	\begin{aligned}
		\sum_{i \in N} U_i(\bmu) &\ge \sum_{i \in N} w_i \\
		U_i(x, \bmu_{-i}) &\le w_i \quad &\forall i \in N, x \in X_i \\
		\sum_{x \in X_i} \mu_i(x) &= 1 \quad &\forall i \in N \\
		0 \le \mu_i(x) &\le 1 \quad &\forall i \in N, x \in X_i
	\end{aligned}	
\end{equation}

As the utility functions are defined by tensors and strategy spaces are finite, the definitions of expected utility simplify into weighted sums. Specifically, the best response function $U_i(x, \bmu_{-i})$ reduces to the vector
\[
	\sum_{\bx_{-i} \in \bX_{-i}} U_i[x,\bx_{-i}] \prod_{j \in N_{-i}} \mu_j[x_j];
\] while the expected utility $U_i(\bmu)$ becomes 
\[
	\sum_{\bx \in \bX} U_i[\bx] \prod_{i \in N} \mu_i[x_i].
\]

As mentioned in \cref{sec:sbab}, algebraic functions can be decomposed into binary operations for which there exist known tight convex relaxations. Specifically, the multilinear program \ref{eq:sturmfels} can be reformulated as a non-convex quadratically-constrained quadratic program, which can then be solved by spatial Branch and Bound solvers.

\todo[inline]{describe our choice and the advantages of each approach}



\begin{figure}
    \centering
    \includegraphics[width=\textwidth]{general.runtimes.pdf}
    \caption{Runtime comparison between the bilinear formulation and the \emph{gambit-logit} algorithm on two-player general-sum matrix games with normally distributed payoffs. Each algorithm was executed 100 times per game size with a 20-second time limit. Mean runtimes of successful runs are overlaid as thick lines. Although the bilinear program exhibits lower median runtimes (shown as darkest plot regions) compared to \emph{gambit-logit} with default settings, its runtime has a multimodal distribution with a significantly higher standard deviation (5.15 s vs. 0.22 s for $30\times30$ games) leading to frequent timeouts. The \emph{gambit-logit} program always finished successfully, while the success rate of the bilinear formulation was steadily decreasing, completing 89\% of the 30-action games and only 59\% of 40-action games. Mean runtimes on 20-action games were: 0.81 s (bilinear) and 0.18 s (\emph{gambit-logit}).}
    \label{fig:general.runtimes}
\end{figure}

\begin{figure}
    \centering
    \includegraphics[width=\textwidth]{testfun.runtimes.pdf}
    \caption{Runtime comparison between the bilinear formulation and the \emph{gambit-logit} algorithm on two-player general-sum matrix games. Each algorithm was executed 100 times per game size with a 20-second time limit. Solid lines represent mean runtimes. Payoff matrices were generated by evaluating the Michalewicz function for Player 1 and McCormick function for Player 2 on a Cartesian product of randomly sampled actions uniformly distributed over $[0, \pi]$. This test should provide a more realistic benchmark for the Double Oracle master-problem compared to normally-distributed random matrices.}
    \label{fig:testfun.runtimes}
\end{figure}

\begin{figure}
    \centering
    \includegraphics[width=\textwidth]{zero.runtimes.pdf}
    \caption{Runtime comparison of the bilinear program, linear program, and the \emph{gambit-logit} algorithm for solving two-player zero-sum matrix games with normally distributed payoffs. Each algorithm was executed 100 times per game size with a 20-second time limit. Thick lines indicate mean runtimes of successful executions. For 20-action games, mean runtimes were: 0.6 ms (bilinear program), 0.02 ms (linear program), and 159.7 ms (\emph{gambit-logit}).}
    \label{fig:zero.runtimes}
\end{figure}
